\documentclass{article}

% typesetting and text macros
\usepackage[letterpaper, textheight=9.0in, textwidth=5.2in]{geometry}
\usepackage{setspace}
\setstretch{1.08}
\usepackage{fancyhdr}
\setlength{\topmargin}{0.0in}
\setlength{\headheight}{3.0ex}\addtolength{\topmargin}{-0.5\headheight} % make room for headers
\setlength{\headsep}{3.0ex}\addtolength{\topmargin}{-0.5\headsep}
\pagestyle{fancy}
%\renewcommand{\headrulewidth}{0.0pt}
\fancyhf{}
\fancyhead[R]{\sffamily\thepage}
\sloppy\sloppypar\raggedbottom\frenchspacing

\usepackage{amsmath}
\DeclareMathOperator{\ee}{e}
\DeclareMathOperator{\ii}{i}
\DeclareMathOperator*{\argmin}{argmin}
\DeclareMathOperator{\loss}{loss}
\DeclareMathOperator{\rank}{rank}

\fancyhead[L]{\sffamily team / Ragged, robust, and heteroskedastic data factorizations}
\title{\bfseries A learned set of continuous eigenfunctions, as a model for heterogeneous and ragged data, by way of robust heteroskedastic factorizations in function space}
\author{Hogg \and others}
\date{April 2026}

\begin{document}

\maketitle

\paragraph{Abstract:}
The original dimensionality reduction method is PCA, which represents a rectangular data matrix with a rank-$K$ approximation, optimizing a mean-squared loss.
There are extensions and improvements, but most of them require rectangular data (though some permit missing values).
In many astrophysical contexts, data are truly ragged; for example,
think of \textsl{LSST} observations of a star or NASA \textsl{SPHEREx} observations of a pixel on the sky:
Each star or pixel gets a different number of observations, at different locations in time and in wavelength.
Here we provide a dimensionality reduction that learns a low-dimensional basis of continuous functions, constrained by heterogeneous and ragged data, and a projection of the ragged data onto the basis.
The method interpolates and extrapolates spectra, answers counterfactual questions about missing data, and identifies outlier and anomalous data points.

\bigskip\noindent
First there was principal components analysis (PCA) \cite{pca},
then there was heteroskedastic matrix factorization (HMF) \cite{hmf},
and then there was Robust HMF \cite{rhmf}.
In all of these, the idea is that the user has \emph{rectangular data}:
There are $N$ data vectors (lists really) $\vec{y}_n$, each of which is an ordered list of $d$ features, making up a $N\times d$ data matrix $Y$.
We will refer to a single element of $Y$ as $y_{nj}$ with $1\leq n\leq N$ and $1\leq j\leq d$.

In PCA, each of these $N\times d$ data points $y_{nj}$ is treated as identically valuable;
in HMF, each data point comes with its own weight $w_{nj}$ (and, indeed, some of those weights can be zero);
in robust HMF, each data point gets a weight adjustment depending on how well the low-rank structure that I am about to describe fits that data point.
Importantly, HMF and robust HMF permit there to be missing data, because data points can be given zero weights.

In all three cases, the fundamental model is a low-rank linear representation, or a matrix factorization.
That is, each model can be written as
\begin{align}
    y_{nj} &= \sum_{k=1}^K a_{nk}\,g_{kj} + \text{residual}_{nj} ~,
\end{align}
where the $a_{nk}$ are coefficients ($K$ for each data vector $n$ and the $g_{kj}$ are elements of eigenvectors or eigenlists, which can be thought of as comprising $K$ $d$-vectors $\vec{g}_k$.
Some might write this equation as
\begin{align}
    Y &= A\,G + \text{residuals} ~,
\end{align}
where $A$ and $G$ are $N\times K$ and $K\times d$ matrices of coefficients and eigenvectors,
and now the residuals make up a $N\times d$ array.
The only differences, in this context, between PCA, HMF, and robust HMF, is what loss function of the residuals is used to set the coefficients $A$ and eigenvectors $G$.
That is, each of these methods can be written in the form
\begin{align}
    \hat{A}, \hat{G} &= \argmin_{A, G} \loss(Y - A\,G) \quad\text{s. t.}~\rank(A\,G) = K ~,
\end{align}
for some loss (different for each method).

In many astronomical contexts, we don't get rectangular data.
For example, in \textsl{LSST} \cite{lsst}, each source $n$ on the sky will have a set of observations taken at different times $t_{nj}$.
No two stars will have the same set of times; indeed different sources will general be observed a different number of times.
In NASA \textsl{SPHEREx} \cite{spherex}, no two pixels on the sky will be observed at the same list of wavelengths; indeed different pixels get different numbers of observations, and each observation $j$ of pixel $n$ is at a different value of wavelength $\lambda_{nj}$ (in general).
For specificity here, we will proceed with the \textsl{SPHEREx} context.

Now the setup is that we have $N$ pixels $n$ on the sky, each of which has been observed a different number $d_n$ of times.
Each pixel $n$ has associated with it a set (not even a list) $\vec{z}_n$, of size $d_n$, of ordered triples:
\begin{align}
    \vec{z}_n &= \left\{ (x_{nj}, y_{nj}, w_{nj}) \right\}_{j=1}^{d_n} ~,
\end{align}
where there are $d_n$ measurements associated with pixel $n$,
each of which is made up of a wavelength (or log-wavelength or other positional parameter) $x_{nj}$, an intensity measurement $y_{nj}$, and a weight (inverse uncertainty variance) $w_{nj}$.
We imagine that each observation $y_{nj}$ is a noisy measurement of a continuous function of the position $x$, such that the equivalent of a low-rank model would be
\begin{align}
    y_{nj} &= \sum_{k=1}^K a_{nk} \sum_{m=-M}^M b_{k,m}\,\ee^{\ii \kappa_m\,x_{nj}} + \text{residual}_{nj} \label{eq:model}\\
    \kappa_m &= \frac{2\pi\,m}{L} ~,
\end{align}
where we have (somewhat arbitrarily) chosen to work in a Fourier-series basis,
the $b_{k,m}$ are complex,
and we will enforce $b_{k,m}=b_{k,-m}^\ast$ to keep everything real in the end.
Along with the rank $K$, we have chosen a real scale $L$ and an integer $M$ to set the domain and the bandwidth of the model.

HOGG: What does this do for you?
It delivers $K$ continuous functions of wavelength---eigenfunctions, if you will---and the best-fit $K$ coefficients for each pixel.
The eigenfunctions $g_k(x)$ are
\begin{align}
    g_k(x) &= \sum_{m=-M}^M b_{k,m}\,\ee^{\ii \kappa_m\,x} ~.
\end{align}
With this definition, the model \eqref{eq:model} can be written
\begin{align}
    y_{nj} &= \sum_{k=1}^K a_{nk}\,g_k(x_{nj}) + \text{residual}_{nj} ~.
\end{align}
This model permits you to predict what you \emph{would have observed} if you had observed pixel $n$ at some wavelength $x_\ast$ at which it has never been previously observed.
It permits you to interpolate and extrapolate the spectrum of a pixel with sparse data.

Optimization proceeds identically to that in robust HMF; we will have an a-step, a b-step, and a w-step, each one analogous to its corresponding step in the robust-HMF work.

HOGG: Get explicit about these steps.

HOGG: Why Fourier series? Include implementation details like \texttt{finufft} \cite{finufft}.

HOGG: Note that you can add a quadratic regularization and thus include a Gaussian-process kernel!

Comment: In PCA, HMF, and RHMF, we deliver the $A$, $G$ outputs from a singular-value decomposition of the low-rank product $A\,G$.
In the continuous-function version, I (Hogg) am not sure this is possible.
Or at least we might have to think about that one.
The magic of sines and cosines might be such that if we orghogonalize (naively) the $B$ matrix, the resulting functions are orthogonal on the domain $L$.

Comment: Non-negative version might be impossible to create?

\bibliographystyle{plain}
\bibliography{ragged}

\end{document}
